\chapter{Line Integrals and Multiple Integrals}\label{ch:b2:multint}

This chapter extends integration from intervals to curves and to
domains of the plane and of space. Line integrals integrate a
\emph{differential form} $P\,\dd x + Q\,\dd y$ along an oriented
arc; double and triple integrals integrate functions over
two- and three-dimensional domains. The two theories meet in the
\emph{Green--Riemann theorem}, the two-dimensional fundamental
theorem of calculus, and the main computational tool throughout is
the \emph{change of variables formula}, whose distortion factor is
the absolute Jacobian determinant.

\section{Line integrals}

\begin{definition}[Differential form; line integral]\label{def:b2:multint:lineint}
Let $U \subseteq \R^2$ be open. A \emph{differential form} of
degree $1$ and class $\mathcal{C}^0$ on $U$ is an expression
$\omega = P\,\dd x + Q\,\dd y$ with $P, Q \colon U \to \R$
continuous --- formally, a continuous map from $U$ into the dual
of $\R^2$, $\omega(M) = P(M)\,e_1^* + Q(M)\,e_2^*$. For a
$\mathcal{C}^1$ arc $\gamma \colon [a, b] \to U$, $\gamma(t) =
(x(t), y(t))$, the \emph{line integral} of $\omega$ along $\gamma$
is
\[
\int_\gamma \omega
= \int_a^b \Bigl(P(\gamma(t))\,x'(t)
+ Q(\gamma(t))\,y'(t)\Bigr)\,\dd t .
\]
The definitions extend verbatim to $\R^3$ (forms $P\,\dd x +
Q\,\dd y + R\,\dd z$) and to piecewise $\mathcal{C}^1$ arcs (sum
over the pieces).
\end{definition}

\begin{proposition}[Invariance and orientation]\label{prop:b2:multint:lineinv}
The line integral is unchanged under an increasing $\mathcal{C}^1$
change of parameter, and changes sign under a decreasing one. It
therefore only depends on the \emph{oriented} geometric arc.
\end{proposition}

\begin{proof}
If $\theta \colon [c, d] \to [a, b]$ is a change of parameter and
$\tilde\gamma = \gamma \circ \theta$, then by the chain rule and
the one-variable change of variables $t = \theta(u)$,
\[
\int_{\tilde\gamma}\omega
= \int_c^d \bigl(P(\gamma(\theta(u)))\,x'(\theta(u))
+ Q(\gamma(\theta(u)))\,y'(\theta(u))\bigr)\,\theta'(u)\,\dd u
= \pm\int_a^b \bigl(Px' + Qy'\bigr)(t)\,\dd t ,
\]
with sign $+$ if $\theta$ is increasing ($\theta(c) = a$) and $-$
if decreasing (the bounds swap).
\end{proof}

\begin{example}[Work of a force; circulation]\label{ex:b2:multint:work}
If $F = (P, Q)$ is a force field, $\int_\gamma P\dd x + Q\dd y =
\int_a^b \langle F(\gamma(t)), \gamma'(t)\rangle\,\dd t$ is the
\emph{work} of $F$ along $\gamma$. For $\omega = -y\,\dd x +
x\,\dd y$ along the counterclockwise unit circle $\gamma(t) =
(\cos t, \sin t)$:
\[
\int_\gamma \omega
= \int_0^{2\pi}\bigl((-\sin t)(-\sin t)
+ \cos t\cos t\bigr)\,\dd t = 2\pi ,
\]
twice the enclosed area --- a first hint of Green--Riemann.
\end{example}

\begin{definition}[Exact and closed forms]\label{def:b2:multint:exact}
The form $\omega = P\,\dd x + Q\,\dd y$ of class $\mathcal{C}^0$
is \emph{exact} on $U$ if there is $f \in \mathcal{C}^1(U)$
(a \emph{potential}) with $\omega = \dd f$, i.e.\ $P = f_x$ and $Q
= f_y$. A $\mathcal{C}^1$ form is \emph{closed} if $P_y = Q_x$ on
$U$.
\end{definition}

\begin{theorem}[Fundamental theorem for line
integrals]\label{thm:b2:multint:ftc}
If $\omega = \dd f$ is exact and $\gamma$ is a piecewise
$\mathcal{C}^1$ arc in $U$ from $A$ to $B$, then
\[
\int_\gamma \omega = f(B) - f(A) .
\]
In particular the integral of an exact form along any closed arc
is zero, and every exact $\mathcal{C}^1$ form is closed.
\end{theorem}

\begin{proof}
$\frac{\dd}{\dd t}f(\gamma(t)) = f_x(\gamma(t))x'(t) +
f_y(\gamma(t))y'(t)$ by the chain rule
(\cref{ch:b2:diffcalc}), so the integrand in
\cref{def:b2:multint:lineint} is the derivative of $t \mapsto
f(\gamma(t))$, and the fundamental theorem of calculus gives the
result on each piece; the intermediate values telescope. Closedness
of exact $\mathcal{C}^1$ forms is Schwarz's theorem: $P_y =
f_{xy} = f_{yx} = Q_x$.
\end{proof}

\begin{example}[Closed does not imply exact]\label{ex:b2:multint:angleform}
On $U = \R^2 \setminus \{0\}$, the \emph{angle form}
\[
\omega = \frac{-y\,\dd x + x\,\dd y}{x^2 + y^2}
\]
is closed (direct computation: both $P_y$ and $Q_x$ equal
$\frac{y^2 - x^2}{(x^2+y^2)^2}$), but its integral along the unit
circle is $2\pi \neq 0$ (same computation as
\cref{ex:b2:multint:work}, divided by $1$): $\omega$ is not exact
on $U$. Locally, $\omega = \dd\theta$ for a determination
$\theta$ of the polar angle; the failure is global --- the angle
cannot be defined continuously around the puncture. On domains
without holes the pathology disappears: on a \emph{star-shaped}
open set, every closed $\mathcal{C}^1$ form is exact
(Poincaré's lemma, \cref{exo:b2:multint:8}).
\end{example}

\section{Double integrals}

We take for granted the theory of the Riemann integral in one
variable (Year 1 volume, and \cref{ch:b2:integration}) and sketch
its two-variable version. A function $f$ continuous on a rectangle
$R = [a, b] \times [c, d]$ has a double integral $\iint_R f$,
defined by Riemann sums over grids exactly as in one variable, and
computed by iteration:

\begin{theorem}[Fubini on a rectangle]\label{thm:b2:multint:fubini}
For $f$ continuous on $R = [a,b] \times [c,d]$,
\[
\iint_R f
= \int_a^b \Bigl(\int_c^d f(x, y)\,\dd y\Bigr)\dd x
= \int_c^d \Bigl(\int_a^b f(x, y)\,\dd x\Bigr)\dd y .
\]
\end{theorem}

\begin{proof}
Set $F(x) = \int_c^d f(x, y)\,\dd y$. Uniform continuity of $f$ on
the compact $R$ makes $F$ continuous (dominated estimate:
$\abs{F(x) - F(x')} \leq (d - c)\sup_y\abs{f(x,y) - f(x',y)}$).
Now subdivide $[a,b]$ and $[c,d]$ into $n$ equal parts, giving a
grid of cells $R_{ij}$ of area $\Delta x\,\Delta y$. On each cell,
$\inf_{R_{ij}} f \cdot \Delta x \Delta y \leq
\int_{x_{i-1}}^{x_i}\int_{y_{j-1}}^{y_j} f(x,y)\,\dd y\,\dd x \leq
\sup_{R_{ij}} f \cdot \Delta x \Delta y$ by monotonicity of the
one-variable integral (applied twice). Summing over cells, the
iterated integral $\int_a^b F$ is squeezed between the lower and
upper Riemann sums of the grid; by uniform continuity both sums
converge to the common value defining $\iint_R f$ as $n \to
\infty$. The same argument applies with the roles of $x$ and $y$
exchanged, so both iterated integrals equal $\iint_R f$.
\end{proof}

\begin{definition}[Elementary domains]\label{def:b2:multint:domain}
A domain $D \subseteq \R^2$ is \emph{$y$-elementary} if
\[
D = \{(x, y) : a \leq x \leq b,\
\varphi_1(x) \leq y \leq \varphi_2(x)\}
\]
with $\varphi_1 \leq \varphi_2$ continuous on $[a,b]$
($x$-elementary: symmetrically). For $f$ continuous on a
$y$-elementary $D$,
\[
\iint_D f
= \int_a^b\Bigl(
\int_{\varphi_1(x)}^{\varphi_2(x)} f(x,y)\,\dd y\Bigr)\dd x ,
\]
and one checks (by extending $f$ by an approximation argument, or
by subdividing) that when $D$ is elementary in both directions the
two iterated integrals agree. Domains cut into finitely many
elementary pieces are handled by additivity, and $\operatorname
{Area}(D) = \iint_D 1$.
\end{definition}

\begin{example}\label{ex:b2:multint:triangle}
On the triangle $D = \{0 \leq x \leq 1,\ 0 \leq y \leq x\}$:
\[
\iint_D xy \,\dd x\,\dd y
= \int_0^1 x\Bigl(\int_0^x y\,\dd y\Bigr)\dd x
= \int_0^1 x\cdot\frac{x^2}{2}\,\dd x = \frac18 .
\]
Swapping the order ($x$ from $y$ to $1$): $\int_0^1
y\bigl(\int_y^1 x\,\dd x\bigr)\dd y = \int_0^1
y\,\frac{1 - y^2}{2}\,\dd y = \frac18$ --- same value, different
computation: choosing the order of integration well is half the
craft.
\end{example}

\begin{theorem}[Change of variables]\label{thm:b2:multint:changevar}
Let $\Phi \colon U' \to U$ be a $\mathcal{C}^1$ diffeomorphism
between open sets of $\R^2$, let $K \subseteq U$ be a compact
domain cut into elementary pieces with $K' = \Phi^{-1}(K)$, and
let $f$ be continuous on $K$. Then
\[
\iint_K f(x, y)\,\dd x\,\dd y
= \iint_{K'} f\bigl(\Phi(u, v)\bigr)\,
\abs{\det J_\Phi(u, v)}\,\dd u\,\dd v .
\]
\end{theorem}
\admitted

\begin{remark}
The complete proof --- approximating $\Phi$ by its differential on
a fine grid and controlling the boundary cells --- is long though
not deep; it is done in full in the measure theory of Year 3, as a
consequence of the Lebesgue theory. The heuristic is the picture
already used for surface area: a small square of side $\dd u$ at
$(u, v)$ is mapped, to first order, onto the parallelogram spanned
by $\Phi_u\,\dd u$ and $\Phi_v\,\dd v$, whose area is
$\abs{\det J_\Phi}\,\dd u\,\dd v$
(\cref{lem:b2:surfaces:lagrange}).
\end{remark}

\begin{example}[Polar coordinates]\label{ex:b2:multint:polar}
$\Phi(\rho, \alpha) = (\rho\cos\alpha,\ \rho\sin\alpha)$ has
\[
J_\Phi = \begin{pmatrix}
\cos\alpha & -\rho\sin\alpha\\
\sin\alpha & \rho\cos\alpha
\end{pmatrix},
\qquad \det J_\Phi = \rho ,
\]
so $\dd x\,\dd y = \rho\,\dd\rho\,\dd\alpha$. For the disk $D_R$
of radius $R$:
\[
\iint_{D_R} e^{-(x^2 + y^2)}\,\dd x\,\dd y
= \int_0^{2\pi}\!\!\int_0^R e^{-\rho^2}\rho\,\dd\rho\,\dd\alpha
= \pi\bigl(1 - e^{-R^2}\bigr)
\xrightarrow[R\to\infty]{} \pi .
\]
Comparing with the square $[-R, R]^2$ (which is squeezed between
the disks $D_R$ and $D_{R\sqrt2}$, all integrands positive) gives
$\bigl(\int_{-\infty}^\infty e^{-x^2}\dd x\bigr)^2 = \pi$:
\[
\boxed{\ \int_{-\infty}^{+\infty} e^{-x^2}\,\dd x = \sqrt{\pi}\ }
\]
--- the Gauss integral again, now by its most famous proof
(compare the one-variable derivation in
\cref{ch:b2:integration}).
\end{example}

\section{The Green--Riemann theorem}

\begin{theorem}[Green--Riemann]\label{thm:b2:multint:green}
Let $K \subseteq \R^2$ be a compact domain which is elementary in
both directions (or a finite union of such glued along segments),
with boundary $\partial K$ a piecewise $\mathcal{C}^1$ closed
curve oriented \emph{counterclockwise} (the domain stays on the
left). For $P, Q$ of class $\mathcal{C}^1$ on a neighbourhood of
$K$:
\[
\oint_{\partial K} P\,\dd x + Q\,\dd y
= \iint_K \Bigl(\frac{\partial Q}{\partial x}
- \frac{\partial P}{\partial y}\Bigr)\,\dd x\,\dd y .
\]
\end{theorem}

\begin{proof}
By additivity of both sides under cutting $K$ into pieces (the
line integrals along interior cuts appear twice with opposite
orientations and cancel), it suffices to treat an elementary
domain. We prove $\oint P\,\dd x = -\iint_K P_y$ on a
$y$-elementary domain $D = \{a \leq x \leq b,\ \varphi_1(x) \leq y
\leq \varphi_2(x)\}$; the identity $\oint Q\,\dd y = \iint_K Q_x$
is symmetric ($x$-elementary), and the theorem is their sum.

Compute the double integral by Fubini and the one-variable
fundamental theorem:
\[
\iint_D \frac{\partial P}{\partial y}\,\dd x\,\dd y
= \int_a^b \bigl(P(x, \varphi_2(x)) - P(x, \varphi_1(x))\bigr)
\,\dd x .
\]
Now the boundary of $D$, counterclockwise, consists of: the lower
graph $y = \varphi_1(x)$ traversed left to right, the right
vertical segment $x = b$ (upward), the upper graph $y =
\varphi_2(x)$ traversed \emph{right to left}, the left vertical
segment $x = a$ (downward). Along the vertical segments $x$ is
constant, so they contribute $0$ to $\oint P\,\dd x$; the graphs,
parametrized by $x$, give
\[
\oint_{\partial D} P\,\dd x
= \int_a^b P(x, \varphi_1(x))\,\dd x
- \int_a^b P(x, \varphi_2(x))\,\dd x
= -\iint_D \frac{\partial P}{\partial y}\,\dd x\,\dd y .
\qedhere
\]
\end{proof}

\begin{corollary}[Area by the boundary]\label{cor:b2:multint:area}
Under the hypotheses of \cref{thm:b2:multint:green},
\[
\operatorname{Area}(K)
= \oint_{\partial K} x\,\dd y
= -\oint_{\partial K} y\,\dd x
= \frac12\oint_{\partial K} x\,\dd y - y\,\dd x .
\]
\end{corollary}

\begin{proof}
Apply Green--Riemann to $(P, Q) = (0, x)$, $(-y, 0)$ and
$\frac12(-y, x)$: each time $Q_x - P_y = 1$.
\end{proof}

\begin{example}[Area of the ellipse]\label{ex:b2:multint:ellipsearea}
For $x = a\cos t$, $y = b\sin t$, $t \in [0, 2\pi]$:
\[
\operatorname{Area}
= \frac12\int_0^{2\pi}\bigl(a\cos t \cdot b\cos t
- b\sin t\cdot(-a\sin t)\bigr)\,\dd t
= \frac{ab}{2}\int_0^{2\pi}\dd t = \pi ab .
\]
\end{example}

\begin{remark}
Green--Riemann explains \cref{ex:b2:multint:angleform}: for a
closed form ($Q_x = P_y$), the integral around the boundary of any
domain contained in $U$ vanishes. The angle form fails to be exact
only because the puncture at the origin prevents the disk bounded
by the unit circle from lying inside $U$ --- line integrals of
closed forms detect the holes of the domain. (Pushed further,
this observation becomes de Rham cohomology.)
\end{remark}

\section{Triple integrals}

The theory extends to three variables with no new idea: Fubini
reduces $\iiint$ to three one-variable integrals (either by
\emph{slicing}: $\iiint_K f = \int\bigl(\iint_{K_z}
f\bigr)\dd z$ over the horizontal slices $K_z$, or by
\emph{stacking}: integrating in $z$ first along vertical sticks),
and the change of variables formula holds with the $3 \times 3$
Jacobian.

\begin{example}[Cylindrical and spherical coordinates]\label{ex:b2:multint:spherical}
\emph{Cylindrical} $(x, y, z) = (\rho\cos\alpha, \rho\sin\alpha,
z)$: $\dd x\,\dd y\,\dd z = \rho\,\dd\rho\,\dd\alpha\,\dd z$.
\emph{Spherical} $(x, y, z) = (r\cos\theta\cos\varphi,\
r\sin\theta\cos\varphi,\ r\sin\varphi)$ ($\theta$ longitude,
$\varphi \in [-\frac\pi2, \frac\pi2]$ latitude): expanding the $3
\times 3$ determinant along the last row,
\[
\det J = r^2\cos\varphi ,
\qquad
\dd x\,\dd y\,\dd z
= r^2\cos\varphi\;\dd r\,\dd\theta\,\dd\varphi .
\]
Volume of the ball of radius $R$:
\[
V = \int_0^R\!\!\int_0^{2\pi}\!\!\int_{-\pi/2}^{\pi/2}
r^2\cos\varphi\;\dd\varphi\,\dd\theta\,\dd r
= \frac{R^3}{3}\cdot 2\pi \cdot 2
= \boxed{\frac43\pi R^3} ,
\]
discharging at last the formula admitted in the volume chapters of
the earlier books.
\end{example}

\begin{example}[Volume by slicing: the cone]\label{ex:b2:multint:cone}
A cone of base area $A$ and height $h$ (apex up, base at $z = 0$):
the slice at height $z$ is the base scaled by the factor $(1 -
z/h)$, of area $A(1 - z/h)^2$. Hence
\[
V = \int_0^h A\Bigl(1 - \frac zh\Bigr)^2\dd z = \frac{Ah}{3} :
\]
the one-third of the school formulas, valid for \emph{any} base
shape --- slicing turns it into the integral of a square.
\end{example}

\section{Exercises}

\begin{exercise}[$\star$]\label{exo:b2:multint:1}
Compute $\int_\gamma y^2\,\dd x + x\,\dd y$ along: (a) the segment
from $(0,0)$ to $(1,1)$; (b) the parabola arc $y = x^2$ from
$(0,0)$ to $(1,1)$. Is the form exact?
\end{exercise}

\begin{exercise}[$\star$]\label{exo:b2:multint:2}
Show that $\omega = (2xy + y^3)\,\dd x + (x^2 + 3xy^2 + 1)\,\dd y$
is closed on $\R^2$, find a potential, and compute
$\int_\gamma\omega$ along any arc from $(0, 0)$ to $(1, 2)$.
\end{exercise}

\begin{exercise}[$\star$]\label{exo:b2:multint:3}
Compute $\iint_D (x + y)\,\dd x\,\dd y$ where $D$ is the domain
bounded by $y = x^2$ and $y = x$ ($0 \leq x \leq 1$), in both
orders of integration.
\end{exercise}

\begin{exercise}[$\star\star$]\label{exo:b2:multint:4}
Using polar coordinates, compute $\iint_D \frac{\dd x\,\dd y}{(1 +
x^2 + y^2)^2}$ over the whole plane (as a limit over disks), and
$\iint_{D'} xy\,\dd x\,\dd y$ over the quarter disk $D' = \{x, y
\geq 0,\ x^2 + y^2 \leq 1\}$.
\end{exercise}

\begin{exercise}[$\star\star$]\label{exo:b2:multint:5}
Compute the area enclosed by the astroid $x = \cos^3 t$, $y =
\sin^3 t$, $t \in [0, 2\pi]$, using
\cref{cor:b2:multint:area}. \emph{(Linearize $\sin^2 t\cos^2 t$.)}
\end{exercise}

\begin{exercise}[$\star\star$]\label{exo:b2:multint:6}
Compute the volume of the solid bounded below by the paraboloid $z
= x^2 + y^2$ and above by the plane $z = 1$, by both methods:
stacking (integrate $1 - x^2 - y^2$ over the unit disk, polar
coordinates) and slicing (horizontal slices are disks of radius
$\sqrt z$).
\end{exercise}

\begin{exercise}[$\star\star$]\label{exo:b2:multint:7}
(Gravitational attraction of a ball --- Newton's theorem, special
case) Show that the volume of the spherical shell $a \leq r \leq
b$ is $\frac43\pi(b^3 - a^3)$ and compute $\iiint_{B}
\frac{\dd x\,\dd y\,\dd z}{r}$ over the ball $B$ of radius $R$
($r$ the distance to the origin). \emph{(Spherical
coordinates.)}
\end{exercise}

\begin{exercise}[$\star\star\star$]\label{exo:b2:multint:8}
(Poincaré's lemma, star-shaped case) Let $U$ be star-shaped with
respect to $0$ (i.e.\ $M \in U \Rightarrow [0, M] \subseteq U$)
and $\omega = P\dd x + Q\dd y$ a closed $\mathcal{C}^1$ form on
$U$. Define
\[
f(x, y) = \int_0^1 \bigl(x\,P(tx, ty) + y\,Q(tx, ty)\bigr)\dd t .
\]
Using differentiation under the integral sign
(\cref{ch:b2:integration}) and $P_y = Q_x$, show that $f_x = P$
and $f_y = Q$: every closed form on a star-shaped open set is
exact.
\end{exercise}

\begin{exercise}[$\star\star\star$]\label{exo:b2:multint:9}
(Dirichlet integral by double integration) Justify and exploit
\[
\int_0^\infty\!\!\int_0^\infty e^{-xy}\sin x\;\dd y\,\dd x
\quad\text{vs}\quad
\int_0^\infty\!\!\int_0^\infty e^{-xy}\sin x\;\dd x\,\dd y
\]
on $[0, A] \times [0, \infty)$: show $\int_0^A \frac{\sin
x}{x}\dd x = \frac\pi2 - \int_0^\infty
e^{-Ay}\frac{y\sin A + \cos A}{1 + y^2}\dd y$ and recover
$\int_0^\infty \frac{\sin x}{x}\,\dd x = \frac\pi2$, comparing
with the parameter-integral proof of
\cref{ch:b2:integration}.
\end{exercise}

\begin{exercise}[$\star\star\star$]\label{exo:b2:multint:10}
(Isoperimetric inequality via Wirtinger) Let $\gamma$ be a simple
closed $\mathcal{C}^1$ curve of length $2\pi$, parametrized by arc
length on $[0, 2\pi]$, enclosing area $A$. Using
\cref{cor:b2:multint:area}, Parseval and the Wirtinger inequality
(exercises of \cref{ch:b2:fourier}), prove $A \leq \pi$, with
equality for the circle. \emph{(Normalize $\int_0^{2\pi} x(s)\dd s
= 0$; write $2A = \oint x\,\dd y - y\,\dd x$ and bound $2A \leq
\int (x^2 + y'^2)$ carefully via $2A = \int_0^{2\pi}(xy' - yx')\dd
s$ and $x^2 + y'^2 \geq 2xy'$.)}
\end{exercise}
